% 本文件由 LaTeX 排版 Agent 根据有效 translation.json 自动生成。
% author=Barnabas; work_id=0124; title=巴拿巴书

\chapter{巴拿巴书}

% structure: p0001 source=h1 target=omitted reason=page-title-already-rendered-by-parent

% structure: p0002 source=h2 target=section reason=relative-heading-rank; base=h2

\section{\WorkTitle{第一章 问候之后，作者宣告他将把自己所领受的一部分传授给他的弟兄}}

诸位儿女，在主耶稣基督的名内，愿你们平安！祂在平安中爱了我们。\par

既见天主的义德之果实在你们中间丰盛，我就极其喜乐，超乎言表，为你们幸福而尊贵的精神，因为你们如此有效地领受了那栽植的属灵恩赐。因此，我内心更加喜乐，盼望得救，因为我确实在你们身上看到从富有的爱之主所倾流的圣神。你们那极受渴望的容貌使我因你们而充满惊异。故此，我对此确信，并在我心中深信：自从我在你们中间开始讲话以来，我明白了许多事，因为主在正义之路上陪伴了我。我也因此有最严格的义务爱你们胜过爱自己的灵魂，因为在你们内住着伟大的信德和爱德，而你们盼望祂所应许的永生。因此，我想到，如果我费心把所领受的一部分传授给你们，这对我服事这样精神的人就已是足够的赏报，于是我便赶紧简短地写信给你们，好使你们除了信德之外，也能有完善的知识。主的道理共有三项：生命的希望，生命的开端和完成。因为主借着先知使我们知道已过和现在的事，也赐给我们未来之事的初果；我们看到这些事一件件成就，就应当以更丰富的信德和更高的精神，怀着敬畏亲近祂。我并非作为你们的教师，而是作为你们中的一员，将简要阐述一些事，使你们在当前情况下更加喜乐。\par

% structure: p0005 source=h2 target=section reason=relative-heading-rank; base=h2

\section{\WorkTitle{第二章 犹太人的祭献已被废除}}

因此，既然时日邪恶，撒殚掌管这世界的大权，我们就当留意自己，殷勤查考主的规诫。恐惧与忍耐乃是我们信德的助手；而坚忍和节制是站在我们这边争战的事。当这些在关乎主的事上保持纯洁时，智慧、明达、学识和知识就与他们一同喜乐。因为祂曾借着众先知启示我们，祂不需要祭献、燔祭或供物，如此说道：\Quote{你们献上许多祭物，于我何益？——上主说。燔祭我已饱足，羔羊的脂油、公牛和山羊的血，我都不喜悦。你们来朝见我的时候，谁向你们要求这些呢？不可再践踏我的院宇。即使你们带来细面，也是枉然。馨香是我所憎恶的恶事；你们的月朔和安息日，我不能容忍。}因此，祂废除了这些事，为使我主耶稣基督的新律法——没有受必需之轭的束缚——能有一种属人的奉献。祂又对他们说：\Quote{难道我吩咐你们的祖先，当他们出埃及地的时候，向我献燔祭和牺牲吗？我反倒这样吩咐你们：你们中间不可有人心里怀恨邻人，也不可爱虚假的誓。} \BibleRef{耶 7:22}因此，我们既然有明悟，就当领悟我们父的仁慈旨意；因为祂对我们说话，愿意我们不像他们那样迷失，而应询问如何亲近祂。祂对我们宣告：\Quote{天主所要的祭献是痛悔的精神；上主悦纳的馨香是荣耀造它的心怀。}因此，弟兄们，我们应当殷勤查考关于我们救恩的事，免得那恶者用诡计潜入，将我们从真正的生命中抛出。\par

% structure: p0007 source=h2 target=section reason=relative-heading-rank; base=h2

\section{\WorkTitle{第三章 犹太人的禁食不是真正的禁食，也不为天主所悦纳}}

祂又就这些事对他们说：\Quote{你们为什么像我今日这样禁食，以致你们的声音可以高声喊叫？我所拣选的禁食，不是人谦卑自己的心灵——上主说。即使你们弯颈如环，披上麻布和灰，也不可称为悦纳的禁食。} \BibleRef{依 58:4}祂对我们说：\Quote{看，我所拣选的禁食是这样：——上主说——不是让人谦卑自己的心灵，而是解开一切邪恶的绳索，松开强硬协议的捆锁，释放受压迫者得自由，撕毁一切不义的契约，把你的饼分给饥饿的人，见到赤身露体的给他衣服穿，把无家可归者领到你家中，见到卑微的人不要轻视，也不可回避自己亲人的需要。那时，你的光明将要破晓，你的伤口要迅速愈合；你的公义要在你面前行走，天主的荣耀将环绕你；那时你呼求，天主必应允你；你还在说话时，祂必说：我在这里。倘若你从自己中间除去压迫的链子、伸指妄誓和怨言，慷慨地将你的饼分给饥饿的人，怜恤受屈的心灵。} \BibleRef{依 58:6}为此，弟兄们，祂长久忍耐，预见到祂所预备的子民将真诚地相信祂的爱者。因为祂预先启示了我们这一切，使我们不致像鲁莽接受他们律法的人那样向前冲。\par

% structure: p0009 source=h2 target=section reason=relative-heading-rank; base=h2

\section{\WorkTitle{第四章 假基督临近了：我们因此要避免犹太人的错误}}

因此，我们这些对眼前事件多有探究的人，理应殷勤查考那些能拯救我们的事。让我们彻底避开一切不义的行为，免得它们抓住我们；让我们憎恶现时的错误，好将我们的爱寄于来世；不要放纵我们的灵魂，让它有能力与罪人和恶人同奔，免得我们变得像他们一样。那最后的绊脚石（或危险之源）临近了，关于它正如哈诺客所说，经上记载：\Quote{为此，主缩短了时期和日子，使祂的爱子能急速前来；祂将来到嗣业中。}先知也这样讲：\Quote{十个王国将在地上执政，之后将兴起一个小王，他要将三个王制服于一体。}同样，达尼尔论及同一件事说：\BibleQuote{我看见第四只兽，邪恶而强有力，比地上所有的兽更凶猛，从它身上长出十只角，又从其中生出一只小角，它把三只大角制服于一体。}所以你们应当明白。此外，我作为你们中的一员，且爱你们每一个人和整体胜过自己的灵魂，更恳求你们：现在要谨慎自守，不要像有些人那样增添自己的罪恶，说：\Quote{盟约既是他们的，也是我们的。}但他们最终失去了它，虽然梅瑟已经接受了它。因为经上说：\BibleQuote{梅瑟在山上禁食四十昼夜，从主领受了盟约，即两块用主的手指写的石版。}\BibleRef{出 31:18}但他们转向偶像，就失去了它。因为主这样对梅瑟说：\BibleQuote{梅瑟，快下去，因为你从埃及地领出来的百姓已经犯罪了。}\BibleRef{出 32:7}梅瑟明白了（天主的旨意），就把两块石版从手中扔下；他们的盟约被打破了，为使所爱者耶稣的盟约铭刻在我们心中，基于因信祂而来的望德。现在，我虽不以为师，却出于爱你们之心，想要给你们写许多事，我已经留心不遗漏我本身所有、有助于你们洁净的事。在这末后的时期，我们要认真留意；因为你们信德的所有过去时光于你们毫无益处，除非在这邪恶的时代，我们作为天主的子女，也抵挡来临的危险。为使那恶者无隙可入，让我们逃避一切虚幻，彻底憎恨恶道的行为。不要退隐独居，仿佛已经成义；而要聚集一处，共同探求什么是有益于你们整体福祉的。因为经上说：\BibleQuote{祸哉，那些自视为智者，自认为明达的人！}\BibleRef{依 5:21}让我们具属神的心志；让我们成为天主完美的殿。尽我们所能，默想对天主的敬畏，遵守祂的诫命，好使我们在祂的律例中喜乐。主要按公道审判世界，不偏待人。各人将按自己所行的受报：若他是义人，他的义德将在他前面；若他是恶人，恶的报应在他面前。留心，免得我们安于逸乐，如同蒙召的人，却在罪中沉睡，而恶的王子获得权柄，将我们逐出主的王国。我的弟兄们，更要留意这一点，当你们回想并看到，在以色列中行了如此大的神迹奇事后，他们竟然（最终）被遗弃。让我们警醒，免得我们被发觉应验了经上所记：\BibleQuote{许多人蒙召，但被选者少。}\par

% structure: p0011 source=h2 target=section reason=relative-heading-rank; base=h2

\section{第五章 新盟约，基于基督的苦难，趋向我们的救恩，却趋向犹太人的毁灭}

因为主为此目的忍受了，把自己的肉身交与朽坏，我们得以藉着罪的赦免而圣化，这赦免是藉着他所洒的血完成的。因为圣经上记载关于他，一部分指以色列，一部分指我们；经上说：\BibleQuote{他为我们的过犯被刺伤，为我们的罪孽被压伤：因他受的鞭伤我们得了痊愈。他被牵到宰杀之地，如同羊羔在剪毛人面前无声。}\BibleRef{依 53:5}因此我们应当深深感激主，因为他既使我们知道过去的事，又赐给我们智慧明白现在的事，并使我们对于将来之事不致愚昧无知。如今，经上说：\Quote{网罗并非为鸟不义地设下。}这是说，那明知正义之路却奔向黑暗之路的人，死得公道。再者，我的弟兄们：若主——他是全世界的主，在创世之初天主曾对他说：\BibleQuote{让我们照我们的肖像，按我们的模样造人，}\BibleRef{创 1:26}你们要明白，他是如何忍受了来自人的苦难。先知们从他获得了恩宠，曾预言关于他。而他（既然必须显现于肉身），为能消灭死亡，揭示从死者中复活，忍受了（他所忍受的），为能成就对祖先的许诺，并为自己预备一个新的子民，同时当他住在地上时，显示他在复活人类后也将审判他们。此外，他教导以色列，行了许多伟大的奇迹异事，向他宣讲了真理，并极其喜爱他。但当他拣选自己的宗徒去宣讲他的福音时，他从那些罪人中拣选，他们罪大恶极，为显示他来\Quote{不是召义人，而是召罪人悔改。}然后他显示自己是天主子。因为他若没有降生成人，人怎能因看见他而得救？因为注视那将要消逝的太阳，它是他双手的作品，人的眼睛尚且不能承受它的光芒。因此，天主子降生成人，为的是把他们迫害他先知致死之罪的恶贯满盈。为此缘故，他忍受了。因为天主说：\Quote{他肉身的创伤来自他们；}和\BibleQuote{当我打击牧人时，羊群就必四散。}\BibleRef{匝 13:7}他自愿这样受苦，因为他必须受死在木架上。因为预言关于他的人说：\Quote{求你救我的灵魂脱离刀剑，用钉子钉住我的肉身，因为恶人的集会起来反对我。}他又说：\BibleQuote{看，我将我的背转给鞭打的人，我的颊转给击打的人，我使我的面容像坚硬的岩石。}\BibleRef{依 50:6}\par

% structure: p0013 source=h2 target=section reason=relative-heading-rank; base=h2

\section{第六章 \Person{基督}的苦难与新盟约曾由先知预言}

因此，当祂完成了诫命时，祂说了什么？\BibleQuote{谁与我争论？让他反对我；或者谁与我进入审判？让他接近主的仆人。}\BibleRef{依 50:8}\BibleQuote{祸哉你们，因为你们都要像衣服一样变旧，蛀虫要吃掉你们。}\BibleRef{依 50:9}先知又说：\Quote{既然祂被安放如一块压碎的大石，看，我在熙雍的地基上投放了一块石头，宝贵、特选、房角石、尊贵的。}接着，祂说了什么？\Quote{凡信赖它的，必永远活着。}那么，我们的盼望是基于一块石头吗？绝非如此。而是〔这话〕因祂以能力坚固了祂的肉体〔作为基础〕；因为祂说：\BibleQuote{祂将我安放如坚固的磐石。}\BibleRef{依 50:7}先知又说：\Quote{匠人所弃的石头，已成为屋角的首石。}他又说：\Quote{这是主所造的伟大而奇妙的子日。}我写给你们较为浅白，使你们可以明白。我是你们爱情的渣滓。那么，先知又说了什么？\Quote{邪恶的集会包围了我，他们围绕我如同蜜蜂围绕蜂巢，}\Quote{他们为我的衣服抽签。}因此，由于祂即将在肉身中显现并受苦，祂的受难已预先显示。因为先知反对以色列说：\BibleQuote{祸哉他们的灵魂，因为他们对自己图谋了邪恶的计谋，}\BibleRef{依 3:9}\Quote{说：让我们捆绑正义者，因为他令我们不悦。}\Person{梅瑟}也对他们说：\BibleQuote{看，主天主说：进入那美好土地，就是主向\Person{亚巴郎}、\Person{依撒格}和\Person{雅各伯}所誓许〔赐给〕你们的土地，并承受它，流奶流蜜之地。}\BibleRef{出 33:1}那么，\OriginalTerm{知识}{Knowledge}说了什么？学习吧：\Quote{信赖，}她说，\Quote{那将要肉身显现给你们的——就是\Person{耶稣}。}因为人是受苦之土，因为\Person{亚当}的造化是来自地面的尘土。那么，这话是什么意思：\Quote{进入那美好土地，流奶流蜜之地？}愿我们的主受赞美，祂在我们里面放置了智慧和认识奥秘事的明悟。因为先知说：\Quote{谁明白主的比喻，除非是那智慧、谨慎并爱主的人？}因此，既然通过赦免我们的罪而更新了我们，祂使我们遵行另一种模式，〔祂的目的〕是我们应拥有孩子的灵魂，因为祂藉其圣神重新造了我们。因为圣经论到我们，当祂对圣子说：\BibleQuote{让我们按我们的肖像，按我们的模样造人；让他们管理海里的鱼、天空的飞鸟和地上的走兽。}\BibleRef{创 1:26}主看到这个美好的受造物——人，说：\BibleQuote{你们要生育繁殖，充满大地。}\BibleRef{创 1:28}这些话〔是〕对圣子说的。再者，我要向你们显示，在末期，祂如何为了我们完成了第二次创造。主说：\Quote{看，我要使末后的如同最初的。}对此，先知宣告：\BibleQuote{进入流奶流蜜之地，并治理它。}\BibleRef{出 33:3}因此，看，我们已被重新塑造，正如祂在另一位先知中又说：\BibleQuote{看，主说，我要从这些人——即主的圣神所预见的那些人——身上除去石心，将肉心放在他们里面，}\BibleRef{则 11:19}因为祂要在肉身中显现，并住在我们中间。因为，我的弟兄们，我们心灵的居所是主的圣殿。\BibleRef{弗 2:21}因为主又说：\Quote{我以何物出现在主我的天主面前，并得光荣？}祂说：\Quote{我要在教会中，在我弟兄中间向你宣认；我要在圣者的集会中赞美你。}我们，就是祂领入那美好土地的人。那么，奶和蜜意味着什么？就是，如同婴儿先靠蜜活，然后靠奶，同样我们，因应许的信德和言语而苏醒并存活，将活着治理大地。但祂上面说：\BibleQuote{让他们增多，治理鱼。}\BibleRef{创 1:28}那么，谁能治理走兽、鱼或天空的飞鸟呢？因为我们必须明白，治理意味着权柄，所以人应命令和管理。如果这目前尚未实现，但祂仍应许了我们。何时？当我们自己也变得完美，以致成为主盟约的继承人时。\par

% structure: p0015 source=h2 target=section reason=relative-heading-rank; base=h2

\section{第七章 禁食与被遣走的山羊是基督的预像}

那么，喜乐之子啊，你们要明白：良善的主已将一切事预先指给我们，使我们知道凡事应向谁献上感恩和赞美。因此，如果天主之子——万有之主，审判活人死人的那一位——受难，使他的创伤能赐给我们生命，我们就该相信，天主之子除非为了我们，决不会受难。此外，当他被钉在十字架上时，他们给他喝醋和胆。请听民众的司祭如何预先指示了这事。主写下了他的诫命，命令凡不守斋戒的都应处死，因为他自己也将要为我们的罪献上圣神的器皿为祭，使那在依撒格被献于祭台上时所立的预像得以完全成全。那么，他在先知书中说了什么？\BibleQuote{让他们带着斋戒，吃那为赎他们一切罪过而献的公山羊。}请留心听：\BibleQuote{只有众司祭可以吃内脏，且不用醋洗。}这是为什么呢？因为你们要给我，就是那将为我的新子民的罪献上我血肉的那一位，喝醋和胆；我独自吃，而百姓则禁食哀哭，披麻蒙灰。这些事发生，是为表明他必须为他们受难。那么，那诫命是如何写的？请注意。取两只上好且相似的山羊，献上它们。司祭要取一只作为赎罪的全燔祭。另一只要怎样处理？\BibleQuote{那只是受诅咒的，}他说。请注意耶稣的\OriginalTerm{预像}{type}如今如何显明。\BibleQuote{你们众人要唾弃它，刺穿它，用朱红羊毛绕在它头上，然后把它赶入旷野。}这一切完成后，那担羊的人把它带到旷野，取下它身上的羊毛，放在一棵称为\OriginalTerm{拉基亚}{Rachia}的灌木上，我们在田野中发现这种灌木时，常吃它的果实。唯有这类灌木的果实是甜的。那么，这又是为什么？请留心。\BibleQuote{一只在祭台上，另一只受诅咒；}为什么那受诅咒的还戴上冠冕？因为他们要在那日看见他，身穿朱红袍子直到脚面；他们将说：这不是我们曾经鄙视、刺透、戏弄、钉死的那一位吗？这确实是那曾宣称自己是天主之子的一位。他与他多么相似！为此，他要求山羊要上好且相似，好叫他们看见他那时来临，因山羊的相似而惊异。看哪，这就是将要受难的耶稣的预像。但为什么他们把羊毛放在荆棘中间？这是置于教会眼前的耶稣的预像。他们把羊毛放在荆棘中，使任何想要取走它的人都必须受许多苦，因为荆棘可怕，只有经过受苦才能得到它。同样，他说：\BibleQuote{那些想要看见我并获得我王国的人，必须通过苦难和痛苦才能得到我。} \BibleRef{宗 14:22}\par

% structure: p0017 source=h2 target=section reason=relative-heading-rank; base=h2

\section{第八章 红母牛是基督的预像}

现在，你们认为这是什么预像呢？以色列曾接到命令，让罪大恶极之人献上一头母牛，宰杀焚烧，然后男孩们取灰，放入器皿中，将紫羊毛和牛膝草缠在棍上，这样男孩们逐一洒在百姓身上，使他们从罪中得到洁净。请留意他如何简单地对你说话。这母牛就是耶稣；献祭的罪人就是那些把他带去宰杀的人。但现在这些人不再有罪，不再被视为罪人。洒水的男孩就是那些向我们宣告罪赦和心灵洁净的人。他们共有十二人，对应以色列的十二支派，主赐给他们传福音的权柄。但为何洒水的是三个男孩？对应亚巴郎、依撒格和雅各伯，因为他们在天主面前是伟大的。为何羊毛放在木头上？因为耶稣藉着木头执掌他的王国，这样，那些信他的人将因十字架而永远活着。为何牛膝草与羊毛一起？因为在他的王国里，日子将邪恶污秽，我们却要得救，而身体受苦的人藉着牛膝草的洁净功效得医治。因此，这些事对我们显明，对他们却隐秘，因为他们没有听从主的声音。\par

% structure: p0019 source=h2 target=section reason=relative-heading-rank; base=h2

\section{第九章 割损礼的属灵意义}

祂又论到我们的耳朵，说祂如何割损了它们和我们的心。主在先知中说：\Quote{他们用耳朵的听觉服从了我。}又说：\Quote{远处的人将听见和知道我所行的事。}\BibleRef{依 33:13}又说：\Quote{你们要在心中受割损，这是上主说的。}\BibleRef{耶 4:4}又说：\Quote{以色列，请听，因为这些话是上主你的天主说的。}\BibleRef{耶 7:2}主的圣神再次宣告：\Quote{谁愿意永远活着？愿他藉着听，听我仆人的声音。}又说：\Quote{天哪，请听；地啊，请侧耳，因为天主说了。}\BibleRef{依 1:2}这些是证明。又说：\Quote{这百姓的首领们，请听上主的话。}\BibleRef{依 1:10}又说：\Quote{孩子们，请听在旷野中呼喊者的声音。}因此，祂割损了我们的耳朵，使我们能听祂的话而相信，因为他们所信赖的割损已被废除。因为祂宣告说，割损不是属肉体的，但因为他们被一个恶天使所迷惑而犯了过犯。祂对他们说：\Quote{这些话是上主你们的天主说的}——（在此我发现一条新的诫命）——\Quote{不要在荆棘中播种，但要向上主自行割损。}为什么祂这样说：\Quote{割损你们内心的顽固，不要硬着脖子？}\BibleRef{申 10:16}又说：\Quote{看，上主说，所有民族都在肉体上未受割损，但这民却在心上未受割损。}\BibleRef{耶 9:25}但你们会说：\Quote{是的，这民确实受了割损，作为印记。}然而，每个叙利亚人、阿拉伯人、以及所有偶像的司祭不也是如此吗？难道他们也在祂的盟约之内？是的，连埃及人也施行割损。孩子们，请学习这一切丰盛的事：那首先吩咐割损的亚巴郎，因着在神视中预见耶稣，便行了此礼，领受了三个字母的奥秘。因为圣经说：\Quote{亚巴郎就给他家中的三百一十八名男子行了割损。}那么，他在这事上得了什么知识呢？先学习十八，然后是三百。十和八是这样表示的：十为\OriginalTerm{十}{\GreekText{Ι}}，八为\OriginalTerm{八}{\GreekText{Η}}。你便有了耶稣名字的首字。又因为十字架要藉字母\OriginalTerm{十字架}{\GreekText{Τ}}表达我们救赎的恩宠，他也说\Quote{三百。}因此，他用两个字母表示耶稣，用一个字母表示十字架。那将祂道理之植入恩赐放在我们内的，知道这事。无人从我这里获得比这更卓越的知识，但我知道你们是堪当的。\par

% structure: p0021 source=h2 target=section reason=relative-heading-rank; base=h2

\section{第十章 梅瑟关于各种食物之诫命的精神意义}

现在，梅瑟为何说：\BibleQuote{你们不可吃猪，也不可吃鹰、隼、乌鸦，以及任何没有鳞的鱼？}他心中［如此］包含了三项教义。此外，主在申命纪中对他们说：\BibleQuote{我要在这百姓中设立我的条例。}\BibleRef{申 4:1}难道不是天主的命令叫他们不可吃［这些］吗？是的，但梅瑟是带着属灵的参照而说的。因此他提到猪，意思就是：\Quote{你们不可与像猪的人结交。}因为当这些人生活在享乐中时，他们忘记了主；但当他陷入困乏时，他们又承认主。同样，猪在吃饱时，不认它的主人；但饥饿时它叫唤，得到食物后又安静了。他说：\BibleQuote{你们也不可吃}他说\BibleQuote{鹰、隼、鸢、乌鸦。}他意思是：\Quote{你们不可与这样的人结交：那些不知靠劳苦和汗水为自己谋食，却以不义夺取他人之物，虽然外表纯朴，却伺机掠夺别人的人。}因此这些鸟，闲坐无事，寻思如何吞吃他人的肉，以邪恶证明自己是［对所有人］的害物。他说：\BibleQuote{你们不可吃}他说\BibleQuote{七鳃鳗、章鱼或乌贼。}他意思是：\Quote{你们不可与那些至终不敬虔、被判定死罪的人结交或仿效他们。}如同那些上述受诅咒的鱼，在深水中漂流，不像其他鱼那样游泳，却栖息在底部的泥中。此外，他说：\BibleQuote{你们不可}他说\BibleQuote{吃野兔。}何故？\Quote{你们不可作败坏男童的人，也不可像这样的人。}因为野兔年复一年增加其受孕的部位；它活多少年，就有多少这样的部位。此外，他说：\BibleQuote{你们不可}他说\BibleQuote{吃鬣狗。}他意思是：\Quote{你们不可作奸淫者、败坏者，也不可像这样的人。}何故？因为那动物每年改变性别，时而雄性，时而雌性。此外，他正确地憎恶了鼬鼠。因为他意思是：\Quote{你们不可像那些我们听说因不洁而用口行恶的人；也不可与那些用口行不义的污秽妇人结交。因为这动物是用口受孕的。}于是梅瑟颁布了关于肉食的三项教义，具有属灵意义；但他们却按肉欲来领受，好像他仅仅是在说实际的肉食。然而，达味领会了这三项教义的知识，也如此说：\BibleQuote{有福的人，没有依从恶人的计谋，}如同那些鱼在黑暗中向深海行去；\BibleQuote{也没有站罪人的道路，}如同那些自称敬畏主、却像猪一样走迷的人；\BibleQuote{也没有坐亵慢者的座位，}如同那些伺机捕食的鸟。要全面而坚定地把握这属灵的知识。但梅瑟更进一步说：\BibleQuote{你们可以吃一切分蹄和反刍的动物。}这是什么意思？［反刍的动物象征］那领受食物后，认出那供养他的恩主，且因他而满足，明显欢喜的人。梅瑟说得很好，尊重了诫命。那么他是什么意思？我们应当与那些敬畏主的人结交，那些在心中默想所领受的诫命的人，那些讲论并遵守主的判断的人，那些知道默想是喜乐的工作，并反刍主的话的人。又分蹄是什么意思？义人虽然行走在这世界，却期盼那将要来临的圣洁境界。看，梅瑟立法多么完善！但他们如何能明了或领悟这些事呢？我们正确理解了他的诫命，便按主的旨意来解释。为此，他割去了我们的耳朵和心，使我们能明白这些事。\par

% structure: p0023 source=h2 target=section reason=relative-heading-rank; base=h2

\section{第十一章 圣洗与十字架在旧约中的预像}

让我们进一步探究，主是否曾预先描绘了[圣洗之]水与十字架。论到水，确实经上论及以色列人写道：他们不应领受那导致罪得赦免的圣洗，而应为自己另寻一个。先知因此宣告说：\Quote{诸天哪，要因此惊奇，大地也要战栗，因为这百姓行了两件大恶：他们离弃了我这活泉，为自己凿出破裂的水池。我的圣山熙雍岂是荒凉的磐石？因为你们必像鸟的雏儿，巢被挪去时便飞走。} \BibleRef{依 16:1} 先知又说：\Quote{我必在你前面行，削平山岳，打破铜门，砍断铁闩；我要将秘密、隐藏、看不见的珍宝赐给你，叫你知道是我，上主天主。} \BibleRef{依 45:2} 以及\Quote{他必住在高峻的磐石洞穴中。} 再者，关于圣子祂说了什么？\Quote{祂的水是可靠的；你要看见君王在祂的荣美中，你的心要默想主的畏惧。} \BibleRef{依 33:16} 又在另一位先知中说：\Quote{行这些事的人必像一棵树栽在水流之旁，按时结果子；叶子也不枯干，凡他所作的尽都顺利。恶人却不是这样，乃像糠秕被风吹散。因此当审判的时候，恶人必站立不住；罪人在义人的会中也是如此。因为主知道义人的道路，恶人的道路却必灭亡。} 请注意祂如何同时描绘了水和十字架。因为这些话意指：有福的是那些寄望于十字架而进入水中的人；因为祂说，他们必按时得赏报；随后祂说：我必报答他们。但现在祂说：\Quote{他们的叶子也不枯干。} 这是说，凡你们口中出于信德和爱德所出的每一句话，都必引领多人悔改和盼望。又有一位先知说：\Quote{雅各伯的土地必超乎万邦之上。} \BibleRef{索 3:19} 这是指祂圣神的器皿，祂要光荣它。再者，祂说了什么？\Quote{有一条河从右边流出来，从它生出美好的树木；凡吃它的果子的人必永远活着。} \BibleRef{厄 47:12} 这是说，我们确实满身罪恶和不洁而进入水中，但出来时心中怀着果实，在灵魂中怀着对天主的敬畏和对耶稣的信赖。\Quote{凡吃这些果子的人必永远活着，} 这是说：祂宣告，凡听见你们说话而相信的人，必永远活着。\par

% structure: p0025 source=h2 target=section reason=relative-heading-rank; base=h2

\section{第十二章 基督的十字架在旧约中的多次预告}

照样，祂在另一位先知中指向基督的十字架，那先知说：\Quote{这些事何时成就？主说：当一棵树被弯下，又再次立起，当血从木头流出。}这里你又得到一个关于十字架和那将要被钉者的暗示。此外，当以色列被外族人攻击时，梅瑟也讲到这一点。为了提醒他们，当他们受攻击时，是因他们的罪恶而被交付于死亡，圣神在梅瑟心中说话，要他做一个十字架和那将要在其上受苦者的形像；因为除非他们信靠祂，他们将永远被战胜。于是梅瑟将一件武器放在另一件之上，立于山丘中央，站在其上，比众人更高，他伸出双手，这样以色列又获得了胜利。但当他放下双手时，他们又被消灭。为什么呢？好让他们知道，除非他们信靠祂，他们不能得救。而在另一位先知中，祂宣告：\BibleQuote{我整天向一个不信的、违抗我正义之道的民族伸出双手。}\BibleRef{依 65:2}梅瑟又做了耶稣的一个预象，表明祂必须受苦，并且祂将成为生命的源泉，而他们却以为祂在十字架上被毁灭了（当时以色列正衰败）。因为由于厄娃借着蛇犯了罪，主使各种蛇咬他们，使他们死亡，\BibleRef{户 21:6}好让他们明白，因着他们的过犯，他们被交于死亡的困境。此外，梅瑟命令：\BibleQuote{你们不可为自己制造任何雕刻或铸成的神像，}为的是揭示耶稣的预象。于是梅瑟制作了一条铜蛇，将它放在一根杆子上，并宣告召集民众。当他们聚集时，他们恳求梅瑟为他们献祭，并为他们祈求痊愈。梅瑟对他们说：\BibleQuote{你们中间若有人被咬，让他来到挂在杆子上的蛇那里；让他盼望并相信，即使它已死，但它能赐给他生命，他立刻就会复原。}\BibleRef{户 21:9}他们就照做了。在此你们也看到耶稣的光荣；因为万有都是在祂内并归于祂。\BibleRef{哥 1:16}再者，梅瑟对农的儿子若苏厄说了什么，当他作为先知给他这个名时，惟独为此，使全体百姓听到：圣父关于祂的圣子耶稣的一切事，都要向农的儿子启示？这个名既已给他，当他派他去窥探那地时，他说：\BibleQuote{你要拿一本书在手中，写下主所宣告的：在末日，天主子要从根本上剪除阿玛肋克全家。}\BibleRef{出 17:14}看哪，再次：那已显现的耶稣，既借着预象，也借肉身显现，\BibleRef{弟前 3:16}祂不是人子，而是天主子。因此，既然他们要说基督是达味之子，他畏惧并明白恶者的错误，他说：\BibleQuote{主对我主说：你坐在我右边，直到我使你的仇敌作你的脚凳。}又，依撒意亚这样说：\BibleQuote{主对我主基督说：我紧握祂的右手，使列邦在祂面前降服；我要打碎君王的力量。}\BibleRef{依 45:1}看，达味如何称祂为主和天主子。\par

% structure: p0027 source=h2 target=section reason=relative-heading-rank; base=h2

\section{基督徒，而非犹太人，是盟约的继承人}

但让我们看看，这民族是继承人，还是先前的民族？盟约属于我们还是属于他们？现在请听圣经关于这民族说了什么。依撒格为他的妻子黎贝加祈祷，因为她不能生育；她就怀了孕。\BibleRef{创 25:21}此外，黎贝加出去求问上主；上主对她说：\BibleQuote{两个国家在你腹中，两个民族在你肚里；这民族要胜过那民族，年长的要服事年幼的。}\BibleRef{创 25:23}你们应当明白谁是依撒格，谁是黎贝加，以及祂关于哪些人宣告这民族将大于那民族。在另一个预言中，雅各伯更清楚地对他儿子若瑟说：\BibleQuote{看，主没有剥夺我见你的面；把你的儿子们带到我这里，好让我祝福他们。}他就带了默纳协和厄弗辣因来，希望默纳协被祝福，因为他是长子。为此若瑟领他到父亲雅各伯的右边。但雅各伯在神视中看见了将来要兴起的民族的预象。圣经说什么？雅各伯改变了他双手的方向，把右手按在次子厄弗辣因的头上，祝福了他。若瑟对雅各伯说：\BibleQuote{请把你的右手移到默纳协的头上，因为他是我的长子。}\BibleRef{创 48:18}雅各伯说：\BibleQuote{我知道，我儿，我知道；但年长的要服事年幼的：然而他也要被祝福。}\BibleRef{创 48:19}你们看见他把手按在谁身上，表明这民族要先来，成为盟约的继承人。如果此外，借着亚巴郎也表明了同样的事，我们的知识便达到圆满。那么，祂对亚巴郎说了什么？\BibleQuote{因为你信了，这就算为你的正义：看，我使你成为那些在未受割礼状态下信主的万民之父。}\par

% structure: p0029 source=h2 target=section reason=relative-heading-rank; base=h2

\section{主赐给了我们那盟约，那是梅瑟领受却又毁坏的}

是的[确实如此]；但让我们查考一下，主是否真的赐下了祂向先祖所起誓要赐给百姓的那盟约。祂确实赐下了，但因了他们的罪，他们不配领受。因为先知宣称：\BibleQuote{梅瑟在西乃山上禁食四十昼夜，为给百姓领受主的盟约。} \BibleRef{出 24:18} 他从主那里 \BibleRef{出 31:18} 领受了两块法版，由主的手的指头藉着圣神所写的。梅瑟领受了之后，就带下去交给百姓。主对梅瑟说：\BibleQuote{梅瑟，梅瑟，快快下去，因为你从埃及地领出来的百姓已经犯了罪。} \BibleRef{出 32:7} 梅瑟明白他们又铸了像，便将法版从他手中扔出去，主的盟约版便摔碎了。梅瑟固然领受了，但他们却证明自己不配。现在你们要学习\Emph{我们}是如何领受的。梅瑟作为仆人 \BibleRef{希 3:5} 领受了它，但主自己为我们受苦，将它赐给了我们，使我们成为嗣业的百姓。但祂显现出来，为的是使他们的不义臻于极处，而使我们藉着祂成为嗣业，领受主耶稣的盟约。主耶稣为此而被预备，好使祂藉着亲自显现，将我们的心（这心已被死亡耗尽，且屈服于错谬的不义）从黑暗中救赎出来，并藉着祂的圣言与我们订立盟约。因为经上记载，圣父将要救赎我们脱离黑暗，命令祂为自己预备圣洁的子民。因此先知宣告：\BibleQuote{我，主，你的天主，凭公义召叫了你，并会握着你的手，加强你，立你作百姓的盟约、万邦的光明，开启盲人的眼目，从锁链中释放被囚的，从黑暗的牢狱中领出坐暗的人。} \BibleRef{依 42:6} 你们因此明白我们从何处被救赎了。先知又说：\Quote{看哪，我已立你作万邦的光明，使你成为救恩，直到地极；救赎你的上主天主这样说。} 先知又说：\BibleQuote{主的神在我身上，因为祂用油膏我，派遣我向谦卑的人传报福音，医治心灵破碎的人，向俘虏宣告释放，向盲人宣告复明，宣布上主恩慈之年和报复之日，安慰一切哀恸的人。} \BibleRef{依 61:1}\par

% structure: p0031 source=h2 target=section reason=relative-heading-rank; base=h2

\section{第十五章 假安息日与真安息日}

此外，关于安息日，在十诫中[主曾面对面与梅瑟在西乃山上说话]也有记载：\BibleQuote{要用洁净的手和纯洁的心，将主的安息日守为圣日。} \BibleRef{出 20:8} 祂又在别处说：\BibleQuote{如果我的儿子们遵守安息日，我就使我的慈爱停留在他们身上。} \BibleRef{耶 17:24} 安息日在创世之初也被提及[如此]：\Quote{天主利用六天完成了祂手的工程，在第七天完成了，并在那天安息，且圣化了它。} 孩子们，请留意这话的意义：\Quote{祂在六天之内完成了。} 这意味着主将在六千年内完成一切，因为在祂看来，一日如千年。祂亲自作证说：\Quote{看哪，今日将如千年。} 因此，孩子们，在六天内，即六千年内，万物将告完成。\Quote{祂在第七天安息了。} 这意味着：当祂的儿子[再次]来临，毁灭恶人的时期，审判不虔敬的人，改变太阳、月亮和星辰时，祂才真正地在第七日安息。此外，祂说：\Quote{你要用纯洁的手和纯洁的心圣化它。} 因此，如果任何人现在能圣化天主所圣化的日子，除非他凡事心中有纯洁，我们就错了。因此，看哪：只有当我们自己领受了恩许，不再有邪恶，万物皆由主更新，能行正义时，那真正安息的人才会正确地圣化这日。那时，我们先被圣化，才能圣化它。祂又对他们说：\BibleQuote{你们的月朔和安息日，我无法容忍。} \BibleRef{依 1:13} 你们明白祂如何说：你们现在的安息日不蒙我悦纳，而是我所造的那一个[即此]，当我使万物安息时，我将开始第八日，即另一个世界的开始。因此，我们也喜乐地遵守第八日，即耶稣从死者中复活的日子。祂显现后，升到天上。\par

% structure: p0033 source=h2 target=section reason=relative-heading-rank; base=h2

\section{第十六章 天主的属神圣殿}

此外，我还要论及圣殿的事，那些可怜的犹太人多么迷惘、陷于错误，他们不相信天主本身，反而相信圣殿，以为那是天主的住所。因为他们几乎像外邦人一样地在圣殿中敬拜祂。但请学习主在废除圣殿时如何说话：\Quote{谁曾用手心量过诸天，用手掌量过大地？难道不是我吗？}\BibleRef{依 40:12}\Quote{上主这样说：\InnerQuote{天是我的宝座，地是我的脚凳；你们要为我建筑什么样的殿宇？或是我安息的地方在哪里？}}\BibleRef{依 66:1}你们看出他们的指望是虚妄的。此外，祂又说：\Quote{看，那些拆毁这殿的人，他们也要重建它。}这事果然发生了。因为他们去打仗，圣殿被敌人摧毁；如今，他们作为敌人的仆人，将要重建它。再者，也曾启示说，城、圣殿和以色列人民都要被交出。因为圣经说：\Quote{在末日，主要将祂牧场的羊、羊圈和塔楼交付毁灭。}这事正如主所说的一样。让我们查考是否仍有天主的圣殿。有的——就在祂自己宣称要建造并完成的地方。因为经上记载：\Quote{当一周完成时，天主的圣殿要在主的荣耀中、以主的名建造。}\BibleRef{达 9:24}因此我发现圣殿确实存在。请学习它将如何以主的名建造。在我们相信天主之前，我们内心的居所是败坏而软弱的，就像一座人手所造的殿。因为它充满了偶像崇拜，是魔鬼的住所，因为我们做了相反天主（旨意）的事。但你要看见，它将以主的名被建造，使主的圣殿在荣耀中建成。如何建造？请这样学习：我们既获得罪的赦免，并将我们的信赖寄于主的名，我们就成为新受造物，从起初重新被塑造。因此，天主确实在我们内居住。如何？祂的圣言、信德；祂的应许召唤；诫命的智慧；教义的命令；祂亲自在我们内说预言；祂亲自居住在我们内；为我们这些被死亡奴役的人打开了圣殿的门，就是口；并通过赐予我们悔改，将我们引入了不朽的圣殿。那么，凡愿意得救的人，不看人，而看那居住在他内、在他内说话的那一位，惊奇于从未听过他用口说出这样的话，他自己也从未想要听这些话。这就是为主所建的属神圣殿。\par

% structure: p0035 source=h2 target=section reason=relative-heading-rank; base=h2

\section{第十七章 书信第一部分的结论}

就我所能并清楚做到的，我怀有希望，按我的愿望，目前关于你们得救所（要求思考）的事，我一样也没有省略。因为如果我写信给你们论及未来的事，你们不会明白，因为这样的知识隐藏在比喻中。这些事就是这样。\par

% structure: p0037 source=h2 target=section reason=relative-heading-rank; base=h2

\section{第十八章 书信的第二部分。两条道路}

但我们现在要转向另一种知识和教义。有两条道路的教义和权柄，一条属于光明，另一条属于黑暗。但这两条道路之间有很大的区别。因为一条路上站的是天主带来光明的天使，但另一条路上站的是撒殚的使者\BibleRef{格后 12:7}。而祂（即天主）确实是主，直到永远；但他（即撒殚）是不法时期的君王。\par

% structure: p0039 source=h2 target=section reason=relative-heading-rank; base=h2

\section{第十九章 光明的道路}

光之道路如下。若有人渴望前往那指定之处，他必须在事工上热忱。因此，为行走此道而赐予我们的知识如下：你当爱那创造你者；你当荣耀那从死亡中救赎你者。你当心地纯朴，精神富足。你当远离那些行走于死亡之路者。你当憎恨天主所不悦之事；你当憎恨一切伪善。你当不违背主的诫命。你当不自高自大，却要谦逊自持。你当不矜夸己荣。你当不对邻人怀恶谋。你当不让鲁莽侵入你灵魂。你当不犯邪淫；你当不行奸淫；你当不诱拐少年。你当不让天主的话语带着任何不洁从你口中发出。你当在责备他人过犯时不徇情面。你当温良；你当和平。你当畏惧你所听见的言语。你当不对弟兄怀怨。你当不存疑心\BibleRef{雅 1:8}，犹豫事情是否成就。你当不妄称主的名。你当爱邻人胜过自己的灵魂。你当不打胎杀婴；也不在出生后将其毁灭。你当不向儿子或女儿吝啬援助，却要从幼年起教导他们敬畏主。你当不贪恋邻人的财物，也不当贪婪。你当不与傲慢者同流合污，却要与义人和谦卑者为伍。当以善事接纳临于你的试炼。你当不心怀二意，也不当口舌两舌，因为两舌是死亡的陷阱。你当顺服主，并因上下同尊天主肖像的缘故，以谦逊和敬畏顺服其他主人。你当不以苦言命令你那同信一主的仆婢，免得你不尊敬那超越二者的天主，因为祂来召叫人，不是按他们的外貌\BibleRef{弗 6:9}，乃是按圣神所预备的\BibleRef{罗 8:29}。你当在一切事上与邻人分享；你当不称任何物为己有；因为你们既共同分享不朽坏之物，何况那些会朽坏之物呢？你当不口舌急躁，因为口是死亡的陷阱。你当尽可能保持灵魂纯洁。你当不轻易伸手索取，却缩手不给。你当爱那对你宣讲主话语的每一个人，如同爱护你的眼珠。你当日夜思念审判的日子。你当每日寻求圣人的面容，或藉言语省察他们，前去劝勉他们，并思虑如何以言语拯救一个灵魂，或藉你的双手为赎你罪而劳动。你当不迟疑地施舍，也不在施舍时抱怨。\Quote{凡求你的，就给他}，你就必知道那赏报的好施主。你当守护你所领受的，不增不减。你当至终憎恨恶者。你当公义审判。你当不制造分裂，却要使争竞者和好。你当承认你的罪。你当不存邪恶良心而祈祷。这就是光明之道。\par

% structure: p0041 source=h2 target=section reason=relative-heading-rank; base=h2

\section{第二十章 黑暗之道}

但黑暗之道乃是弯曲的，充满咒诅；因为它是带着刑罚的永死之道，其中有毁灭灵魂的事，即：偶像崇拜、狂妄自恃、权力的傲慢、伪善、心怀二意、奸淫、凶杀、掠夺、高傲、过犯、诡诈、恶毒、自满、下毒、行邪术、贪婪、不敬畏天主。[在此道中，]也有那逼迫善人的、憎恨真理的、喜爱谎言的、不知义之赏报的、不持守善事的、不以公义审判照顾寡妇和孤儿的、不警醒敬畏天主[却偏向]邪恶的，温良和忍耐远离他们；那些喜爱虚浮、追求赏报、不怜悯贫乏、不扶助劳累之人的人；那些好诽谤、不认识造他们的主、杀害儿童、毁灭天主手工艺的人；那拒绝贫穷之人、欺压受苦之人、袒护富人、不公义审判穷人、在各方面都悖逆的人。\par

% structure: p0043 source=h2 target=section reason=relative-heading-rank; base=h2

\section{第二十一章 结论}

因此，凡学习了经上所记载的诸多主的训令的人，应当遵行它们。因为遵守这些训令的人，必在天国受到荣耀；但选择其他事物的人，必与他的作为一同灭亡。因此将有复活，因此将有报应。我恳求你们这些首领，若你们愿接受我善意的劝告，在你们中间要有那些你们能施以恩惠的人；不要离弃他们。因为那日子临近了，那日一切将连同那恶者一同灭亡。主临近了，他的赏报也临近了。我再三恳求你们：你们要彼此作良善的立法者，彼此作忠信的建议者；从你们中间除去一切伪善。愿掌管全世界的天主，赐给你们智慧、悟性、明达、认识他的判断，以及忍耐。并要受天主教导，殷勤查问主向你们要求什么，并去实行，使你们在审判之日得以安全。如果你们对善事有任何记忆，请记念我，默想这些事，好让我的渴望和警醒都产生一些善果。我恳求你们，以此作为恩惠请求。当你们还在这美好的器皿中时，不要在任何一件事上失败，而要不停地追求它们，并遵守每一条诫命；因为这些事是值得的。因此我更加恳切地写信给你们，尽我所能，为要鼓励你们。再见，你们这些爱与和平的孩子们。愿荣耀的主和一切恩宠的主与你们的灵魂同在。阿们。\par

\SourceRecord{Translated by Alexander Roberts and James Donaldson.；Ante-Nicene Fathers；Vol. 1.；Edited by Alexander Roberts, James Donaldson, and A. Cleveland Coxe.；Buffalo, NY: Christian Literature Publishing Co.,；1885.；<http://www.newadvent.org/fathers/0124.htm>.}
